% !TeX root = ../main.tex

\section{The scope of the static task in Chapter~\ref{ch:courb}}
\label{apx:errata:static-scope}

Chapter~\ref{ch:courb} reports two tasks. One is sequential: predict the category of the
next place a user visits, from the history of that user's check-ins. The other is static:
classify the category of a place from that place's own representation. This section
records what we established after publication about the second one, and it bears on how
that chapter's category results should be read.

The problem is easiest to see by asking what the static task is actually given. Each place in the
Gowalla data carries a fine-grained venue type, a string such as \texttt{Coffee Shop},
\texttt{Seafood}, or \texttt{Airport}. The place embedding that Chapter~\ref{ch:courb} feeds to the
static task is a lookup table on that string: two places of the same venue type receive the same
vector. The target the task must predict is the top-level category, and \texttt{Coffee Shop} belongs
to Food no matter which coffee shop it is.

The input therefore contains the answer. We confirmed this holds without exception across the five
Gowalla state subsets the collection uses: the venue type takes between 284 and 365 distinct values
per state, and in every state each of those values maps to exactly one of the seven top-level
categories. Not one is ambiguous. A model given the venue type and asked for the category is being
asked to reproduce a lookup it already holds, so the reported accuracy on the static task measures
that mapping rather than any learned semantic inference.

Three qualifications matter, and none of them is a softening.

The first is that the sequential task is unaffected. Its input is a history of visits and
its target is a category the model has not seen for that step, so no equivalent identity
exists between input and label. Every claim Chapter~\ref{ch:courb} makes about the
sequential task, including the improvement that the decomposed representation produces,
stands as published.

The second is that Chapter~\ref{ch:cbic} has a version of the same problem, milder but present, and
saying so is more honest than drawing a line between the two. That chapter builds each place's input
feature from the average of its neighbors' categories and deliberately leaves the place itself out,
which looks like it closes the channel. It does not. The place graph is undirected, so a place
appears in its own neighbors' features, and a single graph convolution averages a place together
with its neighborhood. The place's own label comes back at the first hop, diluted among its
neighbors. The difference between the two chapters is therefore one of degree: an exact lookup in
Chapter~\ref{ch:courb}, a diluted average in Chapter~\ref{ch:cbic}.

The third is that Chapter~\ref{ch:mobiwac} does not inherit the problem, for the reason
that chapter gives on its own terms: it replaces the place-level embedding with a
check-in-level representation, whose input is a single visit, and the identity described
above does not arise.

The statement belongs here rather than in the chapter because Chapter~\ref{ch:courb} is a published,
co-authored article, and this appendix is where this collection records the difference between what
a published text says and what we know now. The
measurement is recorded in the repository released with this work.
